\chapter{Wilderness and Extended Care First Aid}\label{ch:wilderness}

Most of this book assumes professional help is minutes away. In remote areas, help may be hours or days away, and the priorities change: you must stabilize, shelter, and sometimes transport the person yourself. This chapter covers the extra decisions and skills needed when evacuation is delayed.

\section{When Help Is Delayed}

The longer care is delayed, the more conservative the management:

\begin{itemize}[leftmargin=*]
  \item \textbf{Minutes away}: Stabilize, summon help, hand over. Standard first aid applies.
  \item \textbf{Hours away}: Prevent worsening, protect from the environment, and plan shelter, warmth, and water.
  \item \textbf{Days away}: You are now the medical team. Prioritize survival (shelter, warmth, water, food), treat the most life-threatening problems first, and consider evacuation.
\end{itemize}

\textbf{WARNING:} \textbf{The "Golden Hours" rule applies differently in the wilderness.} Decision-making shifts from "what can be done" to "what must be done to get the person out alive." Avoid actions you cannot monitor or that risk making the person worse (e.g., aggressive rewarming you cannot sustain).

\section{Assessment When Help Is Far Away}

Repeat the primary survey (Chapter~\ref{ch:assessment}) at regular intervals and extend it:

\begin{enumerate}[leftmargin=*]
\item Scene safety and stabilization: can the person stay here safely, or must they move?
\item Responsiveness and level of consciousness -- record AVPU every 15 minutes
\item Airway, breathing, circulation -- recheck frequently
\item Head-to-toe examination: look for injuries missed in the initial rush
\item Plan: estimate time to help, weather forecast, daylight remaining, available supplies
\end{enumerate}

If the person can walk, staying active maintains warmth and morale -- but never let pain or injury be ignored. If they cannot walk, plan shelter and evacuation.

\section{Improvised Splints and Slings}

Commercial splints are rarely available in the field. Improvise:

\begin{itemize}[leftmargin=*]
  \item \textbf{Rigid splints}: use trekking poles, tent poles, paddles, straight branches, a rolled sleeping pad, or a backpack frame. Pad with clothing, foam, or socks.
  \item \textbf{Soft splints}: roll a blanket, jacket, or sleeping bag around the limb.
  \item \textbf{Buddy splinting}: tape or tie an injured finger or toe to its healthy neighbor with padding between.
  \item \textbf{Anatomic splinting}: secure an injured arm to the chest with a sling; bind an injured leg to the uninjured leg with padding between.
  \item \textbf{Sling}: fold a triangular bandage or shirt into a sling to support the arm and relieve weight.
\end{itemize}

Secure splints above and below the joint nearest the injury. Check color, warmth, and sensation beyond the splint regularly; loosen if swelling makes it tight. Do NOT straighten a deformed limb if straightening causes more pain -- immobilize in the position found.

\section{Water and Hydration}

Dehydration accelerates heat illness, fatigue, and poor judgment.

\begin{itemize}[leftmargin=*]
  \item Drink regularly before thirst appears; thirst signals early dehydration.
  \item Carry or plan water: roughly 3--4 liters per person per day in moderate conditions, more in heat.
  \item Treat water from natural sources: boil for 1 minute (3 minutes above 2,000 m), use a reliable filter, or use chemical treatment per instructions.
  \item Avoid alcohol and excess caffeine, which worsen dehydration.
  \item If water is short, ration activity and travel to the cooler parts of the day.
\end{itemize}

\section{Field Wound Care}

Wounds cannot be kept fully sterile in the field, but infection risk can be reduced:

\begin{enumerate}[leftmargin=*]
\item Clean hands and use gloves or a clean barrier.
\item Rinse the wound with clean water, ideally potable, to remove debris. Do NOT scrub inside deep wounds.
\item Remove visible loose dirt and foreign material with cleaned tweezers if easy; do not dig for embedded material.
\item Apply antibiotic ointment if available and cover with a clean dressing.
\item Change dressings daily or when soaked or dirty, and clean the wound gently at each change.
\item Watch for infection: increasing pain, redness spreading, warmth, swelling, pus, red streaks, or fever. Any of these mean the person needs professional medical care and, if possible, evacuation.
\end{enumerate}

\textbf{WARNING:} \textbf{Do not close deep or dirty wounds with tape or butterfly closures in the field.} Trapping bacteria inside greatly increases the risk of serious infection. Leave them open and seek medical care.

\section{Blisters and Foot Care}

Foot injuries end trips and can become infected far from help:

\begin{enumerate}[leftmargin=*]
\item At the first "hot spot," stop and cover with tape or a dressing to prevent blister formation
\item If a blister forms, leave small intact blisters alone -- the roof is a natural sterile dressing
\item For large, painful, or high-friction blisters, clean the area, puncture at the edge with a sterilized needle, drain without removing the roof, clean, and dress
\item Keep feet dry: change socks, air feet at breaks, and treat any moisture with drying powder
\item Never ignore numbness, deep pain, or skin discoloration in cold weather -- suspect frostbite (Chapter~\ref{ch:environmental})
\end{enumerate}

\section{Cold-Weather Decision-Making}

Hypothermia management is covered in Chapter~\ref{ch:environmental}. In the field, the critical decision is rewarming:

\begin{itemize}[leftmargin=*]
  \item \textbf{If you can keep the person warm and dry for 24 hours}, gentle field rewarming (warm shelter, dry clothing, warm sweet drinks, skin-to-skin warmth) is appropriate for mild hypothermia.
  \item \textbf{If you cannot sustain warmth or evacuation is long}, prioritize getting the person out. A moderately hypothermic person should not be sent out alone.
  \item \textbf{Frostbite}: a frozen extremity should generally NOT be rewarmed in the field if there is any risk of refreezing -- refreezing causes far more damage than keeping the tissue frozen until definitive care. Warm only if you can keep it warm continuously. Do not rub the area.
  \item Handle a severely hypothermic person with extreme gentleness -- rough handling can trigger fatal heart rhythms.
\end{itemize}

\section{Heat-Weather Decision-Making}

\begin{itemize}[leftmargin=*]
  \item Schedule exertion for cooler morning hours; rest in shade during the hottest part of the day.
  \item Replace salt as well as water when sweating heavily: salty snacks or electrolyte drinks.
  \item Heat stroke (hot, altered mental status) requires immediate, aggressive cooling and urgent evacuation -- even a short delay is dangerous. Cool the person continuously during evacuation.
  \item Do not evacuate a heat-exhaustion victim in the hot part of the day if cool rest and fluids resolve the symptoms.
\end{itemize}

\section{Evacuation Decision-Making}

Ask these questions before deciding to move the person:

\begin{enumerate}[leftmargin=*]
\item Is the person's life in danger where they are, or will they die waiting?
\item Can the person be moved without causing serious harm (suspected spinal injury, open fractures, severe head injury)?
\item How far, how long, and how many people does the move require?
\item Is waiting with a plan (shelter, warmth, supplies, signal) safer than a risky move?
\item Who will go for help, and what exact route and instructions will they carry?
\end{enumerate}

\textbf{WARNING:} \textbf{Never attempt to evacuate a person with suspected spinal injury, severe head injury, or unstable chest injury without proper support and equipment} unless they are in immediate, life-threatening danger where they lie. In those cases, move them as a unit with manual in-line neck stabilization and minimal movement.

\textbf{WARNING:} \textbf{Do not split up a group unless the person going for help is certain of the way back, carries a written description of the injured person's condition and location, and can return with help.} Leaving the injured person alone without a plan is a common fatal error.

\section{Signaling and Rescue}

Make it easy to be found:

\begin{itemize}[leftmargin=*]
  \item \textbf{Three} of anything signals distress: three whistle blasts, three shouts, three fires in a triangle, three flashes of light.
  \item Use a mirror or any shiny object to flash toward search aircraft or distant people.
  \item At night, signal with fire (with a careful fuel plan) or a flashlight.
  \item On open ground, make large visible ground signals (X or SOS) with clothing, rocks, or trampled snow.
  \item Stay near the last known position of the group or vehicle unless the reason to move is compelling -- rescuers search where people were last known.
  \item If a rescue aircraft approaches, lie flat and spread arms and legs like an X, or stand with arms raised in a Y ("yes, land here"); do not confuse the pilot.
\end{itemize}

\quicknote{\textbf{The Buddy System}: On any outdoor trip, someone must always know your route and expected return time, and should raise the alarm if you are overdue. This single measure saves more lives than any gear.}

\section{Review Questions}

\begin{enumerate}[leftmargin=*]
\item How do priorities change when professional help is hours or days away?
\item When is it safe to rewarm a frozen extremity in the field, and when should it be kept frozen?
\item Why should deep or dirty wounds NOT be closed with tape in the field?
\item List four evacuation decision questions you must answer before moving a seriously injured person.
\item What does "three" signal in the wilderness, and why is the buddy system essential?
\end{enumerate}
\index{wilderness first aid}
\index{extended care}
\index{improvised splint}
\index{field wound care}
\index{blisters}
\index{foot care}
\index{evacuation}
\index{signaling for rescue}
\index{buddy system}
\index{water treatment}
